\documentclass[10pt,a4paper]{article}

\usepackage[margin=1.7cm]{geometry}
\usepackage{amsmath, amssymb}
\usepackage{graphicx}
\usepackage{booktabs}
\usepackage{hyperref}
\usepackage[numbers,sort&compress]{natbib}
\usepackage{titlesec}
\usepackage{parskip}

\setlength{\parskip}{3pt plus 1pt}
\titlespacing\section{0pt}{8pt plus 2pt}{4pt plus 2pt}
\titlespacing\subsection{0pt}{6pt plus 2pt}{2pt plus 2pt}

\title{\vspace{-1cm}Sprint 1 Handover Notes\\{\large Config-Driven RL for Just-in-Time Adaptive Interventions}}
\author{William Dennis --- University of Bristol \& National University of Singapore}
\date{July 2026}

\begin{document}
\maketitle
\thispagestyle{empty}

\section{Problem Definition}

Physical inactivity is a leading global health risk. Just-in-Time Adaptive Interventions (JITAIs) use mobile/wearable devices to deliver personalised behaviour-change prompts at moments of opportunity~\cite{klasnja2019heartsteps}. Designing effective JITAI policies requires iteration, but real-world trials are slow and expensive.

We formulate the problem as a Markov Decision Process \(\mathcal{M} = (\mathcal{S}, \mathcal{A}, \mathcal{P}, \mathcal{R}, \gamma)\):
\begin{itemize}
  \item \(\mathcal{S}\): 36-state factored space (step\_bin \(\in\) \{inactive, moderate, active\}, sleep \(\in\) \{good, poor\}, day\_of\_week \(\in\) \{weekday, weekend\}, burden \(\in\) \{low, medium, high\})
  \item \(\mathcal{A}\): 4 actions (idle, movement\_suggestion, goal\_reminder, journal)
  \item \(\mathcal{R} = \alpha \cdot f(\text{step\_bin}) + (1 - \alpha) \cdot g(\text{sleep}) - \lambda \cdot \mathbb{I}[\text{action} \neq \text{idle}]\)
  \item 5 decision points per day, 90-day default episode, 50 seeds per agent
\end{itemize}

\section{Existing Work}

\paragraph{HeartSteps V1/V2} (Klasnja et al., 2019~\cite{klasnja2019heartsteps}; Liao et al., 2019~\cite{klasnja2022heartsteps}) --- First RL-enabled MRT for physical activity. V1 ATE +35 steps (p=0.06). V2: Thompson Sampling with dosage tracking and proxy value function; 78.4\% of participants improved, mean +29.75.

\paragraph{PEARL} (Lee et al., 2025~\cite{lee2025pearl}) --- Largest RL-for-PA RCT to date (n=13,463, 4 arms). RL bandit with COM-B informed 155-nudge bank. RL +296 steps vs control (p=0.0002) at 1 month; sustained +210 at 2 months. Provides the real-world effect-size calibration target.

\paragraph{StepCountJITAI} (Karine \& Marlin, 2024~\cite{karine2024stepcountjitai}) --- OpenAI Gym-style simulation environment for PA interventions, DQN $\approx$3,000 return.

\paragraph{Bandit theory} Russo \& Van Roy (2018) establish TS regret bound \(O(\sqrt{dT\log T})\) for linear contextual bandits. Trella et al. (2022) provide design principles for RL in digital health.

\paragraph{LLM-based simulation} Park et al. (2023) demonstrate generative agent simulations of human behaviour. Wang et al. (2024) use LLMs to simulate patient responses (Patient-\(\Psi\)). Lu et al. (2025) benchmark LLM agents for multi-turn human behaviour simulation.

\paragraph{Data-Driven Simulators \& Burden} Wang et al. (2021)~\cite{wang2021optimizing} developed data-driven behavioral simulators incorporating frequency constraints to explicitly mitigate user over-notification and habituation, maximizing long-term intervention responsiveness.

\paragraph{User Heterogeneity \& Pooling} Tomkins et al. (2021)~\cite{tomkins2021intelligentpooling} introduced \emph{IntelligentPooling} to dynamically balance personalization with group-level data sharing. Similarly, Ghosh et al. (2024)~\cite{ghosh2024rebandit} proposed \emph{reBandit}, utilizing random effects and Empirical Bayes to optimize online hyperparameter adjustment in highly noisy mHealth environments.

\paragraph{Count-Based Bandits} Recent work extending Thompson Sampling to non-linear count data models~\cite{liu2025thompson} (e.g., Poisson and Negative Binomial regressions) demonstrates significant reductions in variance and regret when modeling highly skewed proximal health outcomes compared to standard linear baselines.

\paragraph{Power-Constrained Exploration} Yao et al. (2021)~\cite{yao2021power} established meta-algorithms for probability clipping (\(\pi_{min} \le P(A|S) \le \pi_{max}\)) to protect exploration, ensuring that deploying adaptive personalization policies does not compromise the post-trial statistical power required to test overall intervention efficacy.

\section{Our Work: Config-Driven Environment}

The framework is driven entirely by YAML configuration files validated through Pydantic schemas. No source code changes are required to define new MDPs, agents, or experiments. Core components:

\begin{itemize}
  \item \textbf{State:} Factored state with cyclic advances (day\_of\_week) and rolling window counts (burden tracking over last N steps). Day-boundary sleep routing.
  \item \textbf{Transitions:} Three types --- rule\_based (hand-curated), random (Dirichlet tables, baseline), bootstrap (LLM-generated JSON tables).
  \item \textbf{Rewards:} Safe arithmetic expression parser (AST-based). Variables from state factors or action, with configurable mappings.
  \item \textbf{Agents:} 7 implemented --- fixed, random, epsilon-greedy, decaying epsilon-greedy, Thompson Sampling, UCB, HeartSteps V2. All support contextual variants.
  \item \textbf{Evaluation:} Multi-seed benchmark runner, regression test suite\\ (MVP, sprint1\_random, sprint1\_bootstrap), comparison table formatting.
\end{itemize}

\section{Our Work: Synthetic LLM Patients}

We generate patient transition probability tables by prompting LLMs to sample plausible next-state behaviours:

\textbf{Pipeline:} Prompt generation \(\rightarrow\) concurrent API calling (litellm, 200 workers) \(\rightarrow\) JSON parsing \(\rightarrow\) probability table aggregation.

\begin{itemize}
  \item 22{,}320 prompts per persona (full state-action-timestep factorial, N=10 per cell)
  \item 5 distinct behavioural personas (base, goal\_driven, social\_responder, resistant, stable\_maintainer)
  \item \textbf{Total cost: \$0.27} (DeepSeek V4 Flash via OpenCode Zen)
  \item 42 JSON transition tables (6 within-day + 1 day-boundary per persona/initial model)
  \item Analysis pipeline: parse rate, output distributions, cell consistency, burden interactions, step histograms, transition matrix visualizations (85 PDF figures)
  \item 36-paper literature review across 8 dimensions
\end{itemize}

\section{Initial Results}

HeartSteps V2 achieves \textbf{97--100\% of the optimal DP bound} across all 7 persona/initial-model configurations, dominating bandit baselines by 10--40 percentage points:

\begin{table}[ht]
  \centering
  \small
  \begin{tabular}{lccc}
    \toprule
    Persona & Optimal DP & HeartSteps (\% opt) & Best Bandit (\% opt) \\
    \midrule
    Stable Maintainer & 441.8 & 440.8 (99.8\%) & 402.5 (91.1\%) \\
    Base                  & 379.7 & 379.1 (99.9\%) & 327.8 (86.3\%) \\
    Goal-driven           & 414.1 & 414.5 (100.1\%) & 323.0 (78.0\%) \\
    Social Responder      & 313.5 & 306.8 (97.9\%) & 140.1 (44.7\%) \\
    Initial DeepSeek      & 381.6 & 373.6 (97.9\%) & 297.5 (78.0\%) \\
    Initial GLM 5.2       & 374.6 & 374.0 (99.8\%) & 326.8 (87.2\%) \\
    Resistant             & 214.6 & 163.5 (76.2\%) & 151.7 (70.7\%) \\
    \bottomrule
  \end{tabular}
\end{table}

\textbf{Key findings:} (1) Large persona effects --- Stable Maintainer (441) vs Resistant (215). (2) Epsilon-greedy is the best bandit but still 10--40pp behind HeartSteps. (3) Time-of-day features and step activity levels are most predictive. (4) HeartSteps selects idle $\approx$62\% of the time --- personalisation is partly about knowing when \emph{not} to intervene.

\section{What is Next}

\begin{enumerate}
  \item \textbf{Sensor Foundation Models.} Wearable FM POC (2,754 lines): NormWear (768-dim embeddings), FM-guided action space (15 interventions, 3 archetypes), Thompson Sampling with FM context keys.
  \item \textbf{Other LLMs.} Extend bootstrapping to GLM 5.2, Qwen, Claude. Multi-turn persona dialogues for richer state transitions.
  \item \textbf{Better Prompting.} COM-B behaviour-change taxonomy integration. Structured JSON schema enforcement. Chain-of-thought prompting.
  \item \textbf{Synthetic Clinical Results.} Calibrate transitions to PEARL effect sizes (+200--300 steps). Integrate All of Us (n=59K) and UK Biobank (n=100K) for learned transition models.
  \item \textbf{Deep RL agents.} Extend beyond bandits to DQN and policy-gradient methods for temporal credit assignment. Add power-constrained exploration~\cite{yao2021power} and count-based likelihoods~\cite{liu2025thompson}.
\end{enumerate}

% --- References (not counted toward 2-page limit) ---
\bibliographystyle{unsrtnat}
\bibliography{../design/references}

\end{document}
